\documentclass[../../main-detail.tex]{subfiles}
\begin{document}
\chapter{各種オントロジーの分類木}
ここでは，BFO，DOLCE系のDOLCE+DnS Ultralite (DUL)，GFO，YAMATOの
名前付きクラス間の上位分類を比較できる形で示す．木は主として
\relpath{archive/archive002/upper-ontology/data/}以下のOWLファイルに明記された
\texttt{rdfs:subClassOf}から構成した．OWLの制約式と定義クラスは，最上位分類を読み取るために必要な場合を除いて省略している．
多重継承のクラスは，二つ目以降の出現を破線の枠で示す（凡例は概要章）．

\section{BFO}
出典データは\relpath{archive/archive002/upper-ontology/data/bfo.owl}である．

\begin{figure}[H]
  \centering
  \begin{forest}
    dir tree,
    for tree={scale=0.85}
    [entity
      [continuant $\to$]
      [occurrent $\to$]
    ]
  \end{forest}
  \caption{BFOの最上位分類}
\end{figure}

\begin{figure}[H]
  \centering
  \begin{forest}
    dir tree,
    for tree={scale=0.58}
    [continuant
      [generically dependent continuant]
      [independent continuant
        [immaterial entity
          [continuant fiat boundary
            [one-dimensional continuant fiat boundary]
            [two-dimensional continuant fiat boundary]
            [zero-dimensional continuant fiat boundary]
          ]
          [site]
          [spatial region
            [one-dimensional spatial region]
            [three-dimensional spatial region]
            [two-dimensional spatial region]
            [zero-dimensional spatial region]
          ]
        ]
        [material entity
          [fiat object part]
          [object]
          [object aggregate]
        ]
      ]
      [specifically dependent continuant
        [quality
          [relational quality]
        ]
        [realizable entity
          [disposition
            [function]
          ]
          [role]
        ]
      ]
    ]
  \end{forest}
  \caption{BFOのContinuant以下の分類}
\end{figure}

\begin{figure}[H]
  \centering
  \begin{forest}
    dir tree,
    for tree={scale=0.78}
    [occurrent
      [process
        [history]
        [process profile]
      ]
      [process boundary]
      [spatiotemporal region]
      [temporal region
        [one-dimensional temporal region]
        [zero-dimensional temporal region]
      ]
    ]
  \end{forest}
  \caption{BFOのOccurrent以下の分類}
\end{figure}

\section{DOLCE系：DOLCE + DnS Ultralite}
出典データは\relpath{archive/archive002/upper-ontology/data/dolce-lite.owl}である．ファイル名は
\texttt{dolce-lite.owl}だが，内容はDOLCE+DnS Ultralite 4.2であり，古典的な
DOLCEそのものの分類木ではない．
\begin{figure}[H]
  \centering
  \begin{forest}
    dir tree,
    for tree={scale=0.8}
    [Entity
      [Abstract]
      [Event]
      [Object]
      [Quality]
      [Information entity]
    ]
  \end{forest}
  \caption{DULの最上位分類}
\end{figure}

\begin{figure}[H]
  \centering
  \begin{forest}
    dir tree,
    for tree={scale=0.78}
    [Abstract
      [Formal entity
        [Set]
      ]
      [Region
        [Amount]
        [Physical attribute]
        [Social attribute]
        [Space region]
        [Spatio-temporal region]
        [Time interval]
      ]
    ]
  \end{forest}
  \caption{DULのAbstract以下の分類}
\end{figure}

\begin{figure}[H]
  \centering
  \begin{forest}
    dir tree,
    for tree={scale=0.85}
    [Event
      [Action]
      [Process]
    ]
  \end{forest}
  \caption{DULのEvent以下の分類}
\end{figure}

\begin{figure}[H]
  \centering
  \begin{forest}
    dir tree,
    for tree={scale=0.7}
    [Object
      [Agent
        [Person
          [Natural person]
          [Social person]
        ]
        [Physical agent
          [Natural person, dup]
          [Organism]
        ]
        [Social agent, name=dul-social-agent
          [Collective agent
            [Community]
            [Group]
          ]
          [Organization]
          [Personification]
          [Social person, dup]
        ]
      ]
      [Physical object
        [Physical agent, dup]
        [Physical artifact
          [Designed artifact
            [Designed substance]
          ]
        ]
        [Physical body
          [Biological object
            [Organism, dup]
          ]
          [Chemical object]
          [Substance
            [Functional substance
              [Designed substance, dup]
            ]
          ]
        ]
        [Physical place]
      ]
      [Social object, name=dul-social-object]
    ]
  \end{forest}
  \caption{DULのObject以下の分類（AgentとPhysical object）}
\end{figure}

\begin{figure}[H]
  \centering
  \begin{forest}
    dir tree,
    for tree={scale=0.8}
    [Social object
      [Collection
        [Collective]
        [Configuration]
        [Type collection]
      ]
      [Concept
        [Event type
          [Task]
        ]
        [Local concept]
        [Parameter
          [Unit of measure]
        ]
        [Role]
      ]
      [Description
        [Contract]
        [Design]
        [Diagnosis]
        [Goal]
        [Method]
        [Narrative]
        [Norm]
        [Plan
          [Project]
          [Workflow]
        ]
        [Relation
          [Pattern]
          [Social relation]
        ]
        [Right]
        [Theory]
      ]
      [Information object]
      [Place]
      [Situation
        [Plan execution]
        [Transition]
        [Workflow execution]
        [Time-indexed relation
          [Classification]
          [Parthood]
        ]
      ]
      [Social agent]
    ]
  \end{forest}
  \caption{DULのSocial object以下の分類}
\end{figure}

Qualityには，DUL本体で明示された下位クラスはない．Qualityの値域に相当する抽象物は
Region以下に置かれ，QualityとRegionは\texttt{hasRegion}によって結ばれる．

\begin{figure}[H]
  \centering
  \begin{forest}
    dir tree,
    for tree={scale=0.85}
    [Information entity
      [Information object]
      [Information realization]
    ]
  \end{forest}
  \caption{DULのInformation entity以下の分類}
\end{figure}

\section{GFO}
出典データは\relpath{archive/archive002/upper-ontology/data/gfo-basic.owl}である．このOWLでは
CategoryとIndividualに共通する名前付き上位クラスが明記されていないため，以下では
二つの木として示す．

\begin{figure}[H]
  \centering
  \begin{forest}
    dir tree,
    for tree={scale=0.85}
    [Category
      [Concept]
      [Symbol structure]
      [Universal]
    ]
  \end{forest}
  \caption{GFO BasicのCategory以下の分類}
\end{figure}

\begin{figure}[H]
  \centering
  \begin{forest}
    dir tree,
    for tree={scale=0.75}
    [Individual
      [Abstract]
      [Concrete $\to$]
      [Property $\to$]
      [Relator]
      [Role $\to$]
      [Space time entity $\to$]
    ]
  \end{forest}
  \caption{GFO BasicのIndividual以下の分類}
\end{figure}

\begin{figure}[H]
  \centering
  \begin{forest}
    dir tree,
    for tree={scale=0.68}
    [Concrete
      [Perpetuant]
      [Presential
        [Amount of substrate]
        [Material boundary]
        [Material object]
      ]
      [Processual structure
        [Occurrent
          [Change
            [Continuous change]
            [Discrete change]
          ]
          [Event]
          [History]
        ]
        [Process
          [Continuous process]
          [Discrete process]
          [Processual role]
          [State]
        ]
      ]
    ]
  \end{forest}
  \caption{GFO BasicのConcrete以下の分類}
\end{figure}

\begin{figure}[H]
  \centering
  \begin{forest}
    dir tree,
    for tree={scale=0.78}
    [Property
      [Relational role]
    ]
  \end{forest}
  \qquad
  \begin{forest}
    dir tree,
    for tree={scale=0.78}
    [Role
      [Processual role]
      [Relational role]
      [Social role]
    ]
  \end{forest}
  \caption{GFO BasicのPropertyとRole以下の分類}
\end{figure}

\begin{figure}[H]
  \centering
  \begin{forest}
    dir tree,
    for tree={scale=0.7}
    [Space time entity
      [Space entity
        [Spatial boundary
          [Line]
          [Point]
          [Surface]
        ]
        [Spatial region
          [Topoid]
        ]
      ]
      [Time entity
        [Temporal region
          [Chronoid]
        ]
        [Time boundary
          [Left time boundary]
          [Right time boundary]
        ]
      ]
    ]
  \end{forest}
  \caption{GFO BasicのSpace time entity以下の分類}
\end{figure}

\section{YAMATO}
出典データは\relpath{archive/archive002/upper-ontology/data/yamato.owl}である．このファイルには
\texttt{rdf:ID}の重複が5種類あるため，厳格なRDF/XMLローダーでは読み込みに失敗する．
ここではXMLとして読み込み，1194個のクラス要素を1189個の一意なクラスへ統合した上で，
URIで直接指定された1347本の\texttt{rdfs:subClassOf}辺を抽出した．以下では
\texttt{particular}以下の上位層のみを示し，\texttt{CONTEXT}，
\texttt{ROLE-CONCEPT}，\texttt{ROLE-HOLDER}等からのメタレベルの交差辺は省略する．
ノード名は\verb|\tn{英語}{日本語}|で併記する．日本語はOWL内の
\texttt{rdfs:label}を優先し，日本語ラベルがないクラスには訳語を補った．

\begin{figure}[H]
  \centering
  \begin{forest}
    dir tree,
    for tree={scale=0.85}
    [\tn{particular}{個物}
      [\tn{substrate}{基盤} $\to$]
      [\tn{independent entity}{実在物} $\to$]
      [\tn{dependent entity}{従属的実在物} $\to$]
    ]
  \end{forest}
  \caption{YAMATOのparticular以下の最上位分類}
\end{figure}

\begin{figure}[H]
  \centering
  \begin{forest}
    dir tree,
    for tree={scale=0.72}
    [\tn{substrate}{基盤}
      [\tn{space}{空間}
        [\tn{0D space}{0次元空間}]
        [\tn{1D space}{1次元空間}]
        [\tn{2D space}{2次元空間}]
        [\tn{3D space}{3次元空間}]
      ]
      [\tn{substance}{物質}
        [\tn{hydrogen}{水素}]
        [\tn{iron}{鉄}]
        [\tn{oxygen}{酸素}]
      ]
      [\tn{time}{時間}
        [\tn{absolute time}{絶対時間}
          [\tn{time interval}{時区間}]
        ]
        [\tn{relative time}{相対時間}
          [\tn{time duration}{時間長}]
        ]
        [\tn{time point}{時点}
          [\tn{clock time}{時刻}]
          [\tn{date}{日付}]
        ]
      ]
    ]
  \end{forest}
  \caption{YAMATOのsubstrate以下の分類}
\end{figure}

\begin{figure}[H]
  \centering
  \begin{forest}
    dir tree,
    for tree={scale=0.68}
    [\tn{independent entity}{実在物}
      [\tn{abstract}{抽象物}
        [\tn{number}{数}]
        [\tn{point}{点}]
        [\tn{set}{集合}]
        [\tn{structure}{構造}
          [\tn{figure}{形}]
        ]
        [\tn{truth}{真理値}]
        [\tn{unit}{単位}]
      ]
      [\tn{physical}{具体物}
        [\tn{continuant}{持続物}
          [\tn{integral entity}{統合体}]
          [\tn{non-integral}{非統合体}]
        ]
        [\tn{occurrent}{生起物}
          [\tn{event}{イベント}]
          [\tn{stative}{状態的}
            [\tn{process}{プロセス}]
          ]
        ]
      ]
      [\tn{semi-abstract}{準抽象物}
        [\tn{meaning}{内容}
          [\tn{content}{命題}]
        ]
        [\tn{mind}{精神}]
        [\tn{representation}{表現関連}]
        [\tn{representation form}{表現形態}]
      ]
    ]
  \end{forest}
  \caption{YAMATOのindependent entity以下の上位分類}
\end{figure}

\begin{figure}[H]
  \centering
  \begin{forest}
    dir tree,
    for tree={scale=0.68}
    [\tn{dependent entity}{従属的実在物}
      [\tn{generically dependent}{一般的依存}
        [\tn{generic role}{一般ロール}
          [\tn{food}{食糧}]
          [\tn{fuel}{燃料}]
        ]
        [\tn{quality value}{属性値}
          [\tn{categorical}{カテゴリカル}]
          [\tn{quantity}{量}
            [\tn{qualitative quantity}{定性的量}]
            [\tn{quantitative quantity}{定量的量}]
            [\tn{special}{特殊}]
          ]
        ]
      ]
      [\tn{specifically dependent}{個別的依存}
        [\tn{disposition}{傾向}]
        [\tn{feature}{特徴}]
        [\tn{quality}{質}
          [\tn{generic quality}{属性}]
          [\tn{property}{特性}]
        ]
        [\tn{role}{ロール}
          [\tn{continuant role}{持続物ロール}]
          [\tn{occurrent role}{生起物ロール}]
        ]
      ]
    ]
  \end{forest}
  \caption{YAMATOのdependent entity以下の上位分類}
\end{figure}

\end{document}
